\section{论断登记表、见证的出处与开放条目}
\label{app:registers}

作为依据的\src{observer/hypotheses.tsv}{假设登记表}共有 23 行。每行列明标识符、依赖项、精确陈述、前提、状态、见证所在的位置以及修订条件。依赖项排在引用它的论断之前。可容许历史的范围是陈述本身的一部分。例如，原始 11 位示例无法推广到所有比特串，与它在选定的四个串上全部答对这两件事可以同时成立。

\begin{center}
\begin{tabularx}{\linewidth}{@{}l r X@{}}
\toprule
登记状态 & 数目（共 23 项） & 标识符 \\
\midrule
已证明 & 12 & OBS.01--03, R110.01, R110.05--09, R110.12, LOOP.01, E.01 \\
已反驳 & 4 & R110.02--04, R110.10 \\
在实验范围内得到确认 & 3 & AI.01--03 \\
尚待解决 & 2 & R110.11, AI.04 \\
\texttt{DECLARED\_NOT\_DERIVED} & 1 & E.02 \\
\texttt{UN} & 1 & E.03 \\
\bottomrule
\end{tabularx}
\end{center}

“已证明”和“已反驳”均针对相应行中带全部前提的陈述。各行采用的证据形式有所区别，包括数学证明、可通过计算验证的有限构造和穷举结果。AI.01--03 依次记录助手回答与所示信号之间的关系、实时通道和预先录制的通道在新查询下的差异，以及对既有策略在受测任务族中的修正。这三项来自对助手的独立实验。第~\ref{sec:olmo}~节中冻结 OLMo 的计数有各自的报告和测试集。E.01 证明新增谓词的值未由原理论唯一确定。E.02 记录一项附加前提，E.03 记录具体实例的开放状态。

\paragraph{论断的 36 个地址。}
\src{observer/leaf/leaf.tsv}{条目表}将六个领域与六个描述位置交叉排列。领域为 S（来源）、I（同一性）、M（意义）、T（转移）、A（行动）和 O（观察），描述位置为起源、状态、载体、局部论断、见证和界限。在本版本中，32/36 行的状态为 \texttt{KNOWN}，各自给出文件的 SHA-256 摘要以及可查验的定义、函数或文本段落位置。其余 4/36 行的状态为 \texttt{UNKNOWN}，见证来源字段留空，各项问题仍明确保留：
\begin{itemize}
\item I6 涉及因果历史发生分支时主体的数值同一性。记录后继的约定为各分支分配标识符。
\item M6 涉及对用作充分条件的心物联系给出独立的论证。
\item A6 要求对指定见证（R110.11），在同一个严格二值、空间均匀的 Rule~110 系统内部实现整个旁路通道。
\item O6 涉及反馈发生时给定实例的现象体验谓词取何值（E.03）。
\end{itemize}
条目表中的地址 \texttt{T6} 指转移领域的第六个位置，所引的是有噪声通道的判据。条件定理 T6 属于另一套标识系统。\src{observer/leaf/render_leaf.py}{图像验证程序}检查地址的顺序、哈希值与源文件是否相符、所指符号是否存在、TSV 的规范表示和它能否从 SVG 中逐字节恢复。因此，这幅图可以用作表格的可验证表示。

\paragraph{本版本的出处与状态。}
\src{SHA256SUMS}{发布包根目录的清单}给出 PDF、两份 Lean 日志和源文件的校验和。观察者语料的\src{observer/sha256.txt}{独立清单}记录了 48 个文件。初版的\src{observer/provenance/observer-original.sha256}{历史校验和列表}另行保留，共 36 项。发布包采用了该列表中的 35 个原始文件。内部研究任务说明的条目为记录出处而保留在列表中，说明本身未收入公开发布包。导航和见证地址的变动均反映在本版哈希值中。保存的 JSON 报告包含原始观察记录。对之前运行所用的规约，按其在报告对应版本中的文本字节核对。

\paragraph{后续研究的条件。}
R110.11 要求在同一二值动力系统中，为输入、存储、传递和响应给出可验证的编码，并给出全部四条历史的轨迹。AI.04 要求设置新任务，检验在原始位置不可访问时内部状态的转移，同时设置使用完整历史的阳性对照、可逆替换和擦除。附录~\ref{app:gates}~给出 G4 的数值检验规约。要将转移定位于特定子空间，还需在实验前指定层、干预方式、中介度量、受控的替代解释和容差。单一文本载体上的结果尚不能满足这一条件。

编写本专著时已在 Lean~4.35.0 下完成观察者模块规范导入项的完整检查（附录~\ref{app:lean}）。OLMo 各测试集的状态对应所记录的开发集实验。对于持久可塑性假设，须观察到载体组织结构发生跨回合变化，并在控制外部历史的条件下证实这种变化对后续输入处理具有因果作用。E.02 还需要为选择桥接公设给出独立理由，或从具有明确联系前提的相互竞争的桥接公设推出互不相同的操作性预测。E.03 在本版本中仍为 \texttt{UN}。
